\subsection{Correlación cíclica y conservación del orden de la ruta}
\label{sec:rec27-memoria}

El archivo \eqref{act21:tpk:archivo} conserva la concatenación
ordenada de las muestras. El ejemplo siguiente precisa una información
de orden que dos lectores distintos extraen de una palabra de doce
posiciones, conservando el dominio de cada lectura.

Sea $\tau=(\tau_1,\ldots,\tau_{12})\in\{-1,0,1\}^{12}$ una palabra
ordenada, con índices interpretados módulo $12$. Definimos
\[
\begin{aligned}
 G_a&=\{a,a+4,a+8\}\quad(1\leq a\leq4),\\
 \nu_{120}(\tau)&=\sum_{a=1}^4\operatorname{sgn}
 \left(\sum_{m\in G_a}\tau_m\right),\\
 \nu_{270}(\tau)&=\sum_{m=1}^{12}\tau_m\tau_{m+3}.
\end{aligned}
\]
El histograma marginal es $p_j(\tau)=\#\{m:\tau_m=j\}/12$ para
$j\in\{-1,0,1\}$ y su entropía es
$H(\tau)=-\sum_j p_j(\tau)\log p_j(\tau)$, con $0\log0=0$.

\begin{proposition}[Información de orden distinta de la distribución marginal]
Existen palabras con idéntico histograma y distinta correlación
$\nu_{270}$, incluso manteniendo el mismo valor de $\nu_{120}$.
\end{proposition}
\begin{proof}
Considérense
\[
 \tau^{(a)}=(1,1,0,0,0,0,0,0,0,0,0,0),\qquad
 \tau^{(b)}=(1,0,0,1,0,0,0,0,0,0,0,0).
\]
Ambas tienen $p_1=1/6$, $p_0=5/6$ y $p_{-1}=0$, por lo que sus
entropías marginales coinciden. En cada caso, dos de los conjuntos $G_a$
contienen una unidad y los otros dos sólo ceros: $\nu_{120}=2$.
En $\tau^{(a)}$ ningún par de unidades tiene separación cíclica tres,
mientras que en $\tau^{(b)}$ contribuye exactamente el término $m=1$.
Por tanto $\nu_{270}(\tau^{(a)})=0$ y $\nu_{270}(\tau^{(b)})=1$.
\end{proof}

Esta separación identifica qué información conserva el registro ordenado:
las incidencias entre posiciones, además de las frecuencias de los símbolos.
Un estado reducido que retenga únicamente el histograma asigna la misma
salida a las dos palabras y pierde esa correlación. La entropía marginal,
la incertidumbre entre bases y la entropía de entrelazamiento tienen objetos
de definición distintos; su articulación exige conservar respectivamente
la distribución, la transformación de base y el estado con su partición.
